% Options for packages loaded elsewhere
\PassOptionsToPackage{unicode}{hyperref}
\PassOptionsToPackage{hyphens}{url}
\documentclass[
]{article}
\usepackage{xcolor}
\usepackage[margin=1in]{geometry}
\usepackage{amsmath,amssymb}
\setcounter{secnumdepth}{-\maxdimen} % remove section numbering
\usepackage{iftex}
\ifPDFTeX
  \usepackage[T1]{fontenc}
  \usepackage[utf8]{inputenc}
  \usepackage{textcomp} % provide euro and other symbols
\else % if luatex or xetex
  \usepackage{unicode-math} % this also loads fontspec
  \defaultfontfeatures{Scale=MatchLowercase}
  \defaultfontfeatures[\rmfamily]{Ligatures=TeX,Scale=1}
\fi
\usepackage{lmodern}
\ifPDFTeX\else
  % xetex/luatex font selection
\fi
% Use upquote if available, for straight quotes in verbatim environments
\IfFileExists{upquote.sty}{\usepackage{upquote}}{}
\IfFileExists{microtype.sty}{% use microtype if available
  \usepackage[]{microtype}
  \UseMicrotypeSet[protrusion]{basicmath} % disable protrusion for tt fonts
}{}
\makeatletter
\@ifundefined{KOMAClassName}{% if non-KOMA class
  \IfFileExists{parskip.sty}{%
    \usepackage{parskip}
  }{% else
    \setlength{\parindent}{0pt}
    \setlength{\parskip}{6pt plus 2pt minus 1pt}}
}{% if KOMA class
  \KOMAoptions{parskip=half}}
\makeatother
\usepackage{graphicx}
\makeatletter
\newsavebox\pandoc@box
\newcommand*\pandocbounded[1]{% scales image to fit in text height/width
  \sbox\pandoc@box{#1}%
  \Gscale@div\@tempa{\textheight}{\dimexpr\ht\pandoc@box+\dp\pandoc@box\relax}%
  \Gscale@div\@tempb{\linewidth}{\wd\pandoc@box}%
  \ifdim\@tempb\p@<\@tempa\p@\let\@tempa\@tempb\fi% select the smaller of both
  \ifdim\@tempa\p@<\p@\scalebox{\@tempa}{\usebox\pandoc@box}%
  \else\usebox{\pandoc@box}%
  \fi%
}
% Set default figure placement to htbp
\def\fps@figure{htbp}
\makeatother
\setlength{\emergencystretch}{3em} % prevent overfull lines
\providecommand{\tightlist}{%
  \setlength{\itemsep}{0pt}\setlength{\parskip}{0pt}}
\usepackage{bookmark}
\IfFileExists{xurl.sty}{\usepackage{xurl}}{} % add URL line breaks if available
\urlstyle{same}
\hypersetup{
  pdftitle={The Saint Paul Public School Budget},
  pdfauthor={Nick Arnosti},
  hidelinks,
  pdfcreator={LaTeX via pandoc}}

\title{The Saint Paul Public School Budget}
\author{Nick Arnosti}
\date{}

\begin{document}
\maketitle

Saint Paul just approved an increase in property taxes that will provide
an additional \$37 million per year for public schools. This is
approximately \$1,100 for each student in the district. This made me
wonder, what is the current budget? How does this compare to historical
budgets, and to budgets of other nearby districts? And how is this money
being spent?

\section{Quiz}\label{quiz}

\begin{enumerate}
\def\labelenumi{\arabic{enumi}.}
\tightlist
\item
  How much is SPPS spending per student (total expenditures/total
  enrollment)?
\end{enumerate}

\begin{enumerate}
\def\labelenumi{\alph{enumi})}
\tightlist
\item
  Less than \$10,000
\item
  \$10,000-\$15,000\\
\item
  \$15,000-\$20,000
\item
  \$20,000-\$25,000
\item
  More than \$25,000
\end{enumerate}

\begin{enumerate}
\def\labelenumi{\arabic{enumi}.}
\setcounter{enumi}{1}
\tightlist
\item
  Compared to 20 years ago (2004-2005), are (inflation-adjusted)
  per-student expenditures higher or lower today? By what percentage?
\end{enumerate}

\begin{enumerate}
\def\labelenumi{\alph{enumi})}
\tightlist
\item
  Lower, by approximately 20\%
\item
  Lower, by approximately 10\%\\
\item
  Higher, by approximately 10\%
\item
  Higher, by approximately 20\%
\item
  Higher, by more than 20\%
\end{enumerate}

\begin{enumerate}
\def\labelenumi{\arabic{enumi}.}
\setcounter{enumi}{2}
\tightlist
\item
  Among the following 15 Public School Districts (all in the Twin Cities
  Metro Area), where does Saint Paul rank in terms of per-student
  expenditures?
\end{enumerate}

Anoka-Hennepin; Burnsville-Eagan-Savage; Minneapolis;\\
Mounds View; North Saint Paul-Maplewood-Oakdale; Osseo;\\
Roseville; Rosemount-Apple Valley-Eagan; Saint Paul;\\
South Washington County; White Bear Lake; Edina;\\
Minnetonka; Shakopee; Wayzata

\begin{enumerate}
\def\labelenumi{\arabic{enumi}.}
\setcounter{enumi}{3}
\tightlist
\item
  Saint Paul Public Schools spends money from five funds. Two of these
  are ``Building Construction'' and ``Debt Services'' (which pays off
  debts from past construction projects). Last year, how much did Saint
  Paul spend per student from these two funds combined?
\end{enumerate}

\begin{enumerate}
\def\labelenumi{\alph{enumi})}
\tightlist
\item
  Under \$1,000
\item
  \$1,000 - \$2,000
\item
  \$2,000 - \$3,000
\item
  \$3,000 - \$4,000
\item
  Over \$4,000
\end{enumerate}

\begin{enumerate}
\def\labelenumi{\arabic{enumi}.}
\setcounter{enumi}{4}
\tightlist
\item
  General Fund Expenditures are also categorized. Rank the following
  categories (which I placed in a random order) in terms of total
  expenditure. For bonus points, approximately how much money goes to
  each of these categories?
\end{enumerate}

\begin{itemize}
\tightlist
\item
  Special Education
\item
  Student Transportation
\item
  Pupil Support Services
\item
  Instructional Support Services
\item
  Regular Instruction
\item
  Student Activities and Athletics
\end{itemize}

\begin{enumerate}
\def\labelenumi{\arabic{enumi}.}
\setcounter{enumi}{5}
\tightlist
\item
  The majority of the SPPS budget comes from the General Fund, which
  accounts for just over \$23,000 per student. General Fund Expenditures
  are categorized as either ``Site Based'' (costs for staff and other
  expenses at individual schools) or ``District Based'' (costs for
  district administration, curriculum and professional development,
  human resources and finance, transportation, etc). How many of the
  \$23,000 are District Based spending?
\end{enumerate}

\begin{enumerate}
\def\labelenumi{\alph{enumi})}
\tightlist
\item
  Under \$2,000
\item
  \$2,000 - \$5,000
\item
  \$5,000 - \$8,000
\item
  \$8,000 - \$11,000
\item
  Over \$11,000
\end{enumerate}

\section{Answers}\label{answers}

\begin{enumerate}
\def\labelenumi{\arabic{enumi}.}
\item
  \(e)\) SPPS spends approximately \textbf{\$32,000} per student.
\item
  \(e)\) \$32,000 is about \textbf{42\% higher} than it was in the
  2004-2005 school year (when expenditures were \$14,000 per student,
  equivalent to a little over \$23,000 in 2025 dollars).
\item
  Saint Paul ranks \textbf{first} among these 15 districts in
  per-student expenditures, narrowly beating out Minneapolis. On
  average, the other 14 districts spend \$24,000 per student;
  approximately 3/4 of Saint Paul's costs.
\item
  \(e)\) Construction and debt servicing cost Saint Paul approximately
  \textbf{\$7,000 per student} last year.
\item
  The breakdown is as follows (totals do not add to \$23,000 because
  some categories were omitted from the question).
\end{enumerate}

\begin{itemize}
\tightlist
\item
  Regular Instruction (\$9,000/student, 39.1\% of General Fund)
\item
  Special Education (\$4,400/student, 19.1\% of General Fund)
\item
  Student Transportation (\$1,300/student, 5.7\% of General Fund)
\item
  Pupil Support Services (\$1,200/student, 5.2\% of General Fund)
\item
  Instructional Support Services (\$1,100/student, 4.9\% of General
  Fund)
\item
  Student Activities and Athletics (\$300/student, 1.3\% of General
  Fund)
\end{itemize}

A natural follow-up question is, which expenses count towards each of
these categories. In particular, what are ``Pupil Support Services'' and
``Instructional Support Services''? We'll discuss that more below.

\begin{enumerate}
\def\labelenumi{\arabic{enumi}.}
\setcounter{enumi}{5}
\tightlist
\item
  \(d)\) Approximately \textbf{\$9,400/student} of the general fund
  budget goes to district-based expenses.
\end{enumerate}

To help put that number into context, consider the following chart,
which shows site-based and district-based expenditures per student for
each of the 15 districts mentioned previously. Minneapolis and Saint
Paul are in the upper right corner.

\pandocbounded{\includegraphics[keepaspectratio]{/GeneralFundExpenditures.png}}

Note that district-based expenditures range widely: from \$3,600 to
\$9,400. In other words, Saint Paul spends more than 2.5 times as much
as the lowest-cost district on these expenses (an extra \$5,800 per
student). The variation in site-based expenditures is smaller: these
expenditures range from \$10,000 to \$13,700. I would like to do a
deeper dive into what counts as ``District-Based Expenditures'' and why
Saint Paul's costs are much higher than many other districts, but I
think that will have to wait for another post.

\section{High-Level Conclusions}\label{high-level-conclusions}

The conclusions I draw from this quiz are as follows:

\begin{enumerate}
\def\labelenumi{\arabic{enumi}.}
\tightlist
\item
  There is plenty of money available to educate our students.
\item
  SPPS is spending \emph{a lot} on buildings, which I don't see as
  central to providing a good education.
\item
  A lot of money is going to central district staff and to special
  education. We should make sure that these investments are actually
  helping our students.
\end{enumerate}

Below, I elaborate on each point, and provide more context.

\section{There is plenty of money
available}\label{there-is-plenty-of-money-available}

To put the spending in context, SPPS per-student expenditures are

\begin{itemize}
\tightlist
\item
  Higher than daycare costs.
\item
  Approximately equal to in-state tuition, room and board at the
  University of Minnesota.
\item
  Significantly higher than they have been historically.
\item
  Significantly higher than they are in nearby districts.
\end{itemize}

\textbf{All Education is Expensive}

I include the first two reference points because people often speak
about the high cost of daycare and of higher education. I agree that
these are costly! But in fact, they are no more costly than the rest of
our education system.\footnote{There are reasons we should expect a year
  of public education to cost less than daycare or college. For example,
  daycare has much lower child:staff ratios (4:1 for infants, 7:1 for
  toddlers, and 10:1 for pre-school) than public schools do (my wife
  regularly teaches classes with 30+ students). And tuition, room and
  board includes an apartment and all of your meals! \(\newline\)
  Conversely, it's important to recognize that this is not an
  apples-to-apples comparison. Teachers are paid better than daycare
  providers, and ``tuition, room and board'' is what the \emph{student}
  pays, but not necessarily the actual cost: it ignores additional money
  provided by the state. Nevertheless, I think these are useful
  comparisons to make, because most people believe erroneously that
  daycare and college are much more costly than primary and secondary
  school.}

The key difference is how we pay for them. K-12 education costs are
shared by all of society, and spread out over time (I pay into the
system for my entire adult life, even though each of my children will
attend for 12-13 years). Meanwhile, daycare and college expenses are
paid by individual families in the year that they occur.

\textbf{Misleading Marketing}

I include the third and fourth comparisons (to other years and
districts) because marketing materials often provide the impression that
financial support for schools is eroding, or higher elsewhere. As an
example, consider the referendum handout created by SPPS, which includes
the following figures.

\pandocbounded{\includegraphics[keepaspectratio]{/StateFunding.png}}

\pandocbounded{\includegraphics[keepaspectratio]{/LocalLevy.png}}

I trust that the numbers are accurate, but I feel they are misleading.
The point is, school funding comes from many sources! Focusing on just
the state funding allowance of \$7,500/student and the operating levy of
\$1,200/student misses the vast majority of money that is spent.

\textbf{Declining Enrollment is (only) Part of the Story}

As mentioned above, inflation-adjusted per-student costs are up 42\%
over the past 20 years. Some of this can be explained by a 22\% decrease
in enrollment during that time. If enrollment drops but you have the
same buildings and teachers as before, per-student expenditures rise.

However, declining enrollment isn't the whole story, as total
inflation-adjusted expenditures (not per-student) have \emph{risen} by
11\% over that period. In other words, suppose that enrollment increased
by 8,000 - 9,000 students overnight, back to 2005 levels, and the
district managed to add these students without spending a single penny
more than it currently does (which is of course an absurd assumption).
Even after this miracle, SPPS would be spending 11\% more per student
than it did 20 years ago.

\section{Buildings are Expensive}\label{buildings-are-expensive}

As mentioned above, last year SPPS spent over \$7,000 per enrolled
student on construction and debt servicing. This is approximately 22\%
of total expenditures.\footnote{The district regularly claims to spend
  80\% of its money on salaries, but this is misleading because it means
  80\% of the \emph{general fund}, and completely ignores other funds
  (construction, debt servicing, and food) which are almost entirely
  non-salary expenses.} For that much money, the district could rent an
apartment for 12 months for every two students it supports!

For context, I noticed that spending in this category varies widely
across districts. White Bear Lake is spending over \$10,000 per student
on these expenses, but the following chart shows that 10 of the 14
neighboring districts are spending \$4,100 per student or less.

\pandocbounded{\includegraphics[keepaspectratio]{/Buildings.png}}

It is important to note, the \$7,000 per student number is a significant
increase from the historical norm. Five years ago, this number was
\$2,350 (or about \$2,900 after adjusting for inflation). The increased
spending is because of an ``SPPS Builds'' campaign to majorly renovate
facilities. Thus, one could argue that we should amortize these costs
over many years of benefits. This makes some sense, but only if
construction and debt costs truly go down in a few years.

``The greater the impact to student learning, the greater the need to
complete the project.''

But what does this mean in practice?

Let's look at some of the individual projects underway for SPPS Builds:

Hidden River Middle School is undergoing a renovation that will cost
over \$54 million. The school serves just under 600 students, so this is
over \$90,000 per pupil enrolled. (Of course, this is not an annual
number -- the renovations will presumably benefit many years' worth of
students. But still.)

Bruce Vento Elementary: \$90 Million budget, fewer than 500 students
(over \$180,000 per enrolled student). \$6.9 Million/year budget.

Farnsworth Elementary: \$90 Million budget, fewer than 500 students
(over \$180,000 per enrolled student). Budget allocation to school:
about \$6 million (be careful with this).

Barack Obama Elementary (450 students) and Middle Schools (growing from
zero to a maximum of 450), \$72 Million (after school of
\textasciitilde350 on same site closed). Current annual budget: \$6.4
Million (elementary 4.5 + middle 1.9; presumably middle will grow).

Overall, I am just not convinced that this is a good use of money.

\section{Is the Remaining Money Well
Spent?}\label{is-the-remaining-money-well-spent}

\textbf{Special Education}

Special Education (SPED) expenses include SPED Teachers,
Paraprofessionals, some Psychologists and Social Workers, and SPED
Administrators and Support Staff. These expenses come to \$4,400 for
\emph{every student enrolled in the district} (more than 3 times the
transportation budget). Approximately 19\% of students are enrolled in
special education. This means that we spend approximately \$22,000 on
SPED support for each SPED student (in addition to \$9,000 for their
``regular'' teachers)!

I haven't done a deep dive into how this money is being spent. I believe
some of it actually goes to providing services for SPED students at
private schools. But overall, these numbers seem \emph{very} high to me:
both the number of students enrolled in SPED, and the cost per student
in SPED. I am perfectly willing to spend money (even lots of it) to
educate students. But when the amounts are this large, it's important to
look more closely and ask the question, ``Are we sure that this money is
actually helping students learn?''

\textbf{Instructional Support Services}

What are ``Instructional Support Services''? My understanding (which is
probably incomplete) is that this is primarily money spent on
curriculum, professional development, and IT support.

I suspect that most SPPS professional development has almost no impact
on what happens in the classroom. (It would be interesting to survey
teachers to see if they agree.) At the same time, I believe that you can
train someone to be a better teacher. That is, professional development
\emph{can} be very useful. But making it useful requires more than just
``better'' materials. Rather, it requires time, commitment, and likely
significant structural changes.

For context, my wife has worked in many different settings: a public
school on the Navajo Reservation through Teach for America, a Charter
School in Brooklyn serving entirely students of color, a wealthy private
school in Manhattan, and a Gifted and Talented Public school in Saint
Paul. She feels that the Professional Development from Teach for America
and at the Brookly Charter school were \emph{much, much} more impactful
than anything she has done through Saint Paul Public Schools.

While there are many possible reasons for this, I think the biggest one
is how committed these organizations were to it. In Teach for America,
corps members spent roughly six weeks getting full-time coaching, which
included the opportunity to practice techniques under supervision with
summer school students. The charter in Brooklyn required staff to return
to the building 3-4 weeks before students arrived. Teachers were trained
in specific pedagogical techniques, and asked to practice these
techniques. During the school year, students were let out several hours
early \emph{every Friday afternoon} so that staff could do additional
professional development. Teachers were regularly observed and given
feedback. In summary, these schools asked teachers to spend \textbf{a
lot} of time providing and/or receiving professional development.

It's worth noting that Teach For America corps members are new to
teaching, and turnover at the charter school was much higher than at
SPPS (so most teachers were pretty young). As a result, professional
development is perhaps more necessary for these organizations than it is
for Saint Paul Public Schools.

In summary, if we really want professional development to be useful, we
need to set aside a lot more time for it. This seems likely to be
unpopular with parents (if the time is coming out of the regular school
day) and/or the teachers union (if it is coming out of the summer,
evenings or weekends). Furthermore, veteran teachers are likely more
resistant to learning new techniques than are teachers in their
twenties. Without significant structural changes to how professional
development is done, I am inclined to think a lot of the \$1,100/year
going to instructional support services is not making a big impact in
the classroom.

\section{Data Sources}\label{data-sources}

\end{document}
